% Abstract: one paragraph, under 250 words. Built by the future pass from the
% instruction below. Do not process now.

[Write a single-paragraph abstract of under 250 words. Open with the end-to-end
workflow overview, then move to headline metrics, then close on the external
standards strength. (1) Workflow: state that an autonomous Claude Code Opus 4.7
(1M context) Max agent carried one continuous loop from the instruction
specification in \texttt{papers/VVUQ-02/instructions/output-instruct.md}, through
the generated codebase in \texttt{papers/VVUQ-02/codegen/}, through the full
execution record in \texttt{papers/VVUQ-02/execution/}, into this drafted
manuscript, and toward a future full paper, all on \texttt{kevinkawchak/cancer-automated}.
State the thesis in one sentence: for a single autonomous surgical humanoid the
robotic code assurance process, not code generation, is the substantial and
decision-bearing work, and holding verification, validation, and uncertainty
quantification to a higher standard than the code itself is what will make
upcoming physical AI oncology trial deliverables faster, less expensive, and more
beneficial to patients. (2) Metrics: pull the exact figures from the
\texttt{Key quantitative results to cite} table in \texttt{papers/VVUQ-02/execution/README.md}
and the surface in \texttt{papers/VVUQ-02/execution/03-vvuq/README.md}, namely
172 of 172 automated tests passing with 64 of them in the 10-gate assurance suite,
ten humanoid-specific gates each requiring a verification fraction of exactly 1.0
with validation agreement up to 1.00, relative error as tight as 0.01, and a
coefficient-of-variation bound as tight as 0.05, five decision cases resolving to
ten ACCEPT, three BLOCK mechanisms, and one ESCALATE, all 32 of 32 sweep
iterations clearing every gate at a composite mean near 93.6, the 1790-line
four-entrant tournament, and the 1000-row humanoid sensor stream, with all three
large artifacts reproduced byte-for-byte from seed 20260525. (3) Standards
strength: close by stating that binding every gate to published consensus
standards already used in real life (ASME V\&V 40-2018, NASA-STD-7009A,
IEC 80601-2-77, ISO 13482, ISO/TS 15066, IEC 62304, ISO 14971, UL 4600, and IEEE
7009, resolved at runtime from \texttt{papers/VVUQ-02/inputs/standards/manifest.yaml})
is what makes the assurance argument traceable and defensible to a regulator and
lets an inexpensive autonomous run stand as credible evidence for a hypothetical
2030 platform. Keep the H2-Surgical 1.0 labelled hypothetical and the comparison
simulation-against-simulation. Do not cite in the abstract.]
